% Section - Classes
% Roberto Masocco <roberto.masocco@uniroma2.it>
% April 12, 2026

% ### Classes ###
\section{Classes}
\graphicspath{{figs/classes/}}

% --- class vs struct ---
\begin{frame}{Classes}{\texttt{class} vs \texttt{struct}}
	\justifying
	C++ keeps the \texttt{struct} keyword from C, but gives it new powers: member functions, constructors, and access control.\\
	\bigskip
	It also introduces \texttt{class}. The \textbg{only difference} between the two is the \textbf{default access specifier}:
	\begin{itemize}
		\item members of a \texttt{struct} are \textbg{public} by default;
		\item members of a \texttt{class} are \textbg{private} by default.
	\end{itemize}
	\bigskip
	By convention:
	\begin{itemize}
		\item \texttt{struct} is used for \textbg{plain data aggregates} with no significant behavior;
		\item \texttt{class} is used when \textbg{encapsulation} and \textbg{invariants} matter.
	\end{itemize}
\end{frame}

% --- Access specifiers ---
\begin{frame}{Classes}{Access specifiers}
	\justifying
	Access specifiers control the \textbg{visibility} of class members to the outside world.\\
	\bigskip
	There are three:
	\begin{itemize}
		\item \textbg{\texttt{public}}: accessible from anywhere;
		\item \textbg{\texttt{protected}}: accessible from the class itself and from \textbf{derived classes};
		\item \textbg{\texttt{private}}: accessible \textbf{only} from the class itself.
	\end{itemize}
	\bigskip
	A specifier applies to all members declared \textbf{below it}, until the next specifier.\\
	\bigskip
	Good practice is to expose only what is \textbg{necessary}:
	\begin{itemize}
		\item data members are almost always \texttt{private};
		\item the \textbg{public interface} is kept minimal and stable.
	\end{itemize}
\end{frame}

% --- Constructors and destructors ---
\begin{frame}{Classes}{Constructors and destructors}
	\justifying
	Every class can define \textbg{special member functions} that the compiler calls automatically.\\
	\bigskip
	The \textbg{constructor} is called when an object is \textbf{created}:
	\begin{itemize}
		\item its name is the class name, it has no return type;
		\item it initializes data members and acquires resources;
		\item a class may have \textbg{multiple constructors} with different parameter lists (\textbf{overloading});
		\item if none is defined, the compiler generates a \textbg{default constructor}.
	\end{itemize}
	\bigskip
	The \textbg{destructor} is called when an object is \textbf{destroyed}:
	\begin{itemize}
		\item its name is the class name preceded by \texttt{\textasciitilde};
		\item it releases resources and performs cleanup;
		\item there is \textbg{exactly one} destructor per class, it takes no arguments.
	\end{itemize}
	\begin{alertblock}{Guarantee}
		\justifying
		\textbf{The destructor is always called}, regardless of how the object's scope ends — even in the presence of exceptions.
	\end{alertblock}
\end{frame}

% --- Member functions and this ---
\begin{frame}{Classes}{Member functions and the \texttt{this} pointer}
	\justifying
	Functions declared inside a class are called \textbg{member functions} (or \textit{methods}).\\
	They have direct access to all data members of the class, regardless of access specifiers.\\
	\bigskip
	Inside any non-static member function, the keyword \textbg{\texttt{this}} is an implicit pointer to the \textbf{object on which the function was called}.\\
	\bigskip
	It is rarely needed explicitly, but it is always there:
	\begin{itemize}
		\item it allows disambiguating between a local variable and a member with the same name;
		\item it allows returning a reference to the object itself (\texttt{return *this}).
	\end{itemize}
	\bigskip
	Member functions that do not modify the object should be declared \textbg{\texttt{const}}:
	\begin{itemize}
		\item the compiler then enforces the constraint;
		\item they can be called on \texttt{const} objects or through \texttt{const} references.
	\end{itemize}
\end{frame}

% --- Header and source split ---
\begin{frame}{Classes}{Separating declaration from definition}
	\justifying
	In C++, a class is typically \textbg{declared} in a header file (\texttt{.hpp}) and \textbg{defined} in a source file (\texttt{.cpp}).\\
	\bigskip
	The \textbf{header} contains:
	\begin{itemize}
		\item the class declaration with all member names and types;
		\item inline function definitions (small, frequently called functions).
	\end{itemize}
	\bigskip
	The \textbf{source file} contains:
	\begin{itemize}
		\item the full definitions of all non-inline member functions;
		\item member function names are qualified with \texttt{ClassName::}.
	\end{itemize}
	\bigskip
	\begin{block}{Why split?}
		\justifying
		The header is the \textbg{public contract} of the class — it is what other translation units include.
		Keeping implementation in \texttt{.cpp} files reduces \textbg{compile-time dependencies} and hides internals.
	\end{block}
\end{frame}
